Výstupom z procesu nanopórového sekvenovania je sekvencia nukleotidov (A, T, C, G pre DNA alebo A, U, C, G pre RNA) reprezentujúca genetickú informáciu pôvodnej molekuly \cite{Zheng2023Nanopore}.
Tento výstup je často vo forme textového súboru (napríklad FASTQ formát), ktorý obsahuje sekvencie spolu s kvalitatívnymi skóre pre každú bázu \cite{Zheng2023Nanopore}.
Vzhľadom na technické obmedzenia a variabilitu v procese sekvenovania je bežné, že jednotlivé DNA alebo RNA reťazce sú prečítané niekoľkokrát (tzv. \enquote{coverage} alebo \enquote{depth of coverage}) \cite{Zheng2023Nanopore}.
Toto opakované čítanie zvyšuje spoľahlivosť a presnosť výslednej sekvencie, pretože umožňuje korekciu chýb a identifikáciu variácií v genetickom materiáli \cite{Zheng2023Nanopore,Lopez2019DNA}.
Napriek tomu, že sa reťazce čítajú viackrát, stále môže dôjsť k nepresnostiam v sekvencii v dôsledku šumu v signáli \cite{Zheng2023Nanopore}.
